% ============================================================
% ABSTRACT -- Writer: Ignacio Gimeno
% Every number appears in Section 4. Recheck after Adam's and Pablo's sections
% are merged (US result and selection numbers).
% ============================================================

Universities often rank students from different degrees on a single grade point average (GPA) list to allocate exchange places, grants or admission to master's programmes, which assumes that a grade means the same in every field. We ask how much of a student's position in such a ranking depends on their field of study, and how ranking students within their degree would change who is selected. We analyse three separate datasets: 12,125 mean graduate grades from 89 Spanish universities (2015--16 to 2023--24), the records of 3,703 Portuguese students and course grades from 33 US public universities. Within Spanish universities, 65\% of the variation in grades between fields is a field effect shared by all universities and 28\% is specific to each university. In a supervised regression task with five models tuned on 2018--2021 and tested once on 2022--2023, our main model, a shrunk profile of each university's past deviations, reduces the prediction error from 0.225 to 0.138 grade points relative to the shared field effect; gradient boosting improves on it only slightly (0.133). In a simulated top-20\% selection, 41\% of Nursing graduates are selected against 6\% of Business administration graduates. Subtracting each university's past field deviation from the marks lowers the average distance from equal access from 11.8 to 5.1 percentage points, but cannot remove differences in the spread of grades. Ranking students within their degree reduces the spread of selection rates across degrees from 0.090 to 0.006, while the next-semester grade of those selected changes only from 14.21 to 14.08 out of 20. A single raw-GPA ranking therefore sorts students largely by field, and ranking within the degree corrects this at almost no cost.

\BAMkeywords{grade comparability; GPA ranking; field of study; mixed-effects models; temporal validation; fair selection}
